\section{Introduction}
\label{sec:intro}

Embedding-based retrieval is the first stage of friend recommendation at social-network scale. A graph neural
network embeds every user from its sampled neighbourhood, and nearest-neighbour search over those embeddings proposes
the shortlist of 1{,}000 candidates that later stages rank~\cite{huang2020ebr,ying2018graph,johnson2019billion}. Our
graph has close to a billion users, and the retrieval model is trained contrastively: each user is pulled towards the
friends it went on to add and pushed away from negatives~\cite{yi2019sampling}. The negatives decide much of what such a
model learns, and at this scale they also decide what training costs, because every negative is a user whose
neighbourhood must be sampled and encoded.

Industry practice offers two ways to make negatives harder than random users. The first is logged exposures: users who
were shown a suggestion and did not act on it~\cite{huang2020ebr}. The second is to mine negatives with the current
model~\cite{xiong2021ance}, which means re-embedding candidates as the model changes. On our graph neither works well.
Adding exposure negatives to an otherwise identical training run lowers macro recall@1000 by 0.35 points, with every
ranking metric down, while filling the same slots with random users changes nothing: the harm comes from which users
the exposures are, not from having more negatives (Section~\ref{sec:exposure}). Mining with the current model is expensive
at this scale. From our measured inference rate and step times, refreshing an index of all users once per epoch and
encoding one mined negative per training row would take about four times the GPU time of random training
(Appendix~\ref{app:ance}).

We start instead from a property of the model the production system already trains. In an embedding space learned
with random negatives only, similarity follows the graph: a user's one-hop neighbours average a cosine similarity of
0.70 with the user, three-hop users 0.45 and random users 0.07 (Figure~\ref{fig:locality}). A user's neighbours are even
closer to the user than the friends the user will add next. Two consequences follow. Grouping users by embedding groups
them by graph region, so each user's in-batch negatives come from its own neighbourhood and are hard. And because
grouped users share neighbours, the batch's sampled subgraph shrinks: clustered batches sample 35\% fewer nodes, and a
clustered training step is 12--15\% faster at full scale (Section~\ref{sec:locality}).

Clustered-anchor batching turns this property into a training signal. Every 50 updates, the anchors of the next window
are clustered by their query embeddings and each batch is built from a few clusters. The embeddings come from a cache
that the training loop fills as a by-product, and the clustering runs on the CPU while the GPU trains. On the
production graph, clustered batching raises macro recall@1000 from 34.8\% to 36.4--36.5\% after three epochs, a gain of
1.6--1.7 points, and a second seed reproduces every arm within 0.1 point. The cached version trains on 7\% fewer
GPU-hours than random batching, including the one-time pass that fills the cache (Section~\ref{sec:method}). In a
ten-day online A/B test, the deployed model, trained with this mining and without exposure negatives, increased new
friends made with bidirectional communication by 2.09\%.
\todo{Decide whether NGAME~\cite{dahiya2023ngame} is named here or only in Section~\ref{sec:related}.}

The concentration that makes clustered batches informative can also make them harmful. When a few clusters hold far
more positive rows than the rest, the scheduler fills whole batches, and runs of batches, from one region. The model
then fits that region and forgets the others: on Collab, Citation2 and LiveJournal, uncapped clustered mining ends its
mining phase 8 to 34 points of recall@1000 below random training, and its late batches are drawn from a single cluster
(Section~\ref{sec:public}). A per-cluster row cap bounds the rows any one cluster may contribute to a batch. It removes
the collapse on all three graphs in every seed: capped mining stays within half a point of random at recall@1000, and
raises full-corpus MRR on Citation2 and LiveJournal. Whether a graph needs the cap is visible in how concentrated its
uncapped batches become. Our production graph sits at the safe end, because data preparation keeps at most four
positives per user; there the cap is neutral, and the deployed model runs with it at no cost (Section~\ref{sec:capping}).

Our contributions are:
\begin{itemize}
  \item A measurement on a billion-node social graph: embeddings trained with random negatives are graph-local, which
        makes embedding-clustered batches both harder and cheaper to train on.
  \item Evidence that exposure negatives, the conventional industry source of hard negatives, hurt retrieval on this
        graph, with a slot-matched random control showing that the harm lies in the identity of the exposed users.
  \item Clustered-anchor mining with a cached clustering signal: +1.6 to +1.7 points of macro recall@1000 over random
        batching at full scale, at 7\% lower training cost.
  \item An online A/B test of the deployed model: +2.09\% new friends made with bidirectional communication and
        +1.01\% new friends made.
  \item A failure mode of clustered mining, regional forgetting under concentrated batches, and a per-cluster row cap
        that prevents it on three public graphs, with a batch-concentration statistic that shows when the cap is needed.
\end{itemize}
